% Generated by make_xcomp_definitions.py -- DO NOT EDIT BY HAND.
% Rebuild:  python3 python/make_xcomp_definitions.py --tex
%           cd Text_files && pdflatex xcomp_definitions   (x2)
\documentclass[11pt,a4paper]{article}
\usepackage[margin=2.2cm]{geometry}
\usepackage{amsmath,amssymb,longtable,xcolor,booktabs}
\usepackage[colorlinks=true,linkcolor=black,citecolor=blue,urlcolor=blue]{hyperref}
\title{$X_{\rm comp}$: definitions, dependencies, and the physics behind it}
\author{Compiled for the \texttt{photon\_vs\_electron} project}
\date{2026-09-15}
\begin{document}\maketitle

\section{What $X_{\rm comp}$ is}
At luminosity parity $L_\gamma=L_e$, $X_{\rm comp}$ compares how many ionizations each
agent buys with a given amount of injected energy:
\begin{equation}
  X_{\rm comp}(\alpha_\gamma,\alpha_e)
  = \frac{\displaystyle\int S_\gamma N_{{\rm ion},\gamma}\,dE \big/ \langle E\rangle_\gamma}
         {\displaystyle\int S_e N_{{\rm ion},e}\,dE \big/ \langle E\rangle_e}
  = \frac{\langle N_\gamma\rangle/\langle E\rangle_\gamma}
         {\langle N_e\rangle/\langle E\rangle_e}
  = \frac{W_e(\alpha_e)}{W_\gamma(\alpha_\gamma)},
  \label{eq:xcomp}
\end{equation}
where $W\equiv\langle E\rangle/\langle N_{\rm ion}\rangle$ is the \emph{energy price of one
ionization} --- the $W$-value of radiation physics. $X_{\rm comp}>1$ means photons are the
cheaper ionizing agent per erg.

\paragraph{Everything else cancels.} $n_{\rm H}$, the comoving conversion, the redshift and
the luminosity normalisation all drop out of eq.~\eqref{eq:xcomp} identically, so the map
depends on $\alpha_\gamma$ and $\alpha_e$ \emph{and nothing else}. Because
$\log X_{\rm comp}=\log W_e(\alpha_e)-\log W_\gamma(\alpha_\gamma)$ is additively separable,
the map has no degenerate direction: both axes do real work.

\section{Two repairs to the construction as originally written}
\paragraph{(1) The denominator double-counted the yield.} The integrand was specified as
$S(E)N_{\rm ion}(E)/\langle E_{\rm ion}\rangle$ with $\langle E_{\rm ion}\rangle$ defined as
``the average energy spent per ionization'', i.e.\ $W$ itself. With $S$ number-normalised
that gives $\langle N\rangle/W=\langle N\rangle^{2}/\langle E\rangle$: the yield enters
twice. The form is \emph{dimensionally} correct either way, which is why it does not
announce itself --- but the two readings differ by
$\langle N_\gamma\rangle/\langle N_e\rangle\approx1/106$ and \textbf{invert the conclusion}:
$X_{\rm comp}=2.07$ (photons win) versus $0.0196$ (electrons win $51\times$). Adopted here:
$\langle E_{\rm ion}\rangle=\langle E\rangle$, the mean energy per \emph{primary}, which
makes the integral the ionizations produced per unit energy injected.

\paragraph{(2) The band top is derived, not adopted.} ``The free-streaming regime
($\sim$keV)'' evaluates, using this project's own $\sigma_{\rm photoion}$ over one Hubble
length at $z=10$, to $\tau_{\rm IGM}=1$ at $E=1218.0$~eV. A smooth $(1-e^{-\tau})$
weighting was tried first and \emph{diverges}: it charges the numerator $\langle E\rangle$
for photons that free-stream and ionize nothing while the denominator saturates, so
$W_\gamma$ grows without bound (measured $+224\%$ at $\alpha_\gamma=0.5$ between a $10$~keV
and a $100$~keV integration limit). The sharp band is therefore kept.

\section{The physics: why electrons are the more expensive agent}\label{sec:why}

\paragraph{The floor.} Ionizing hydrogen costs $E_{\rm th}=13.598$~eV and no less, so
$E_{\rm th}/W$ is a true efficiency. At the fiducial indices
($\alpha_\gamma=2.0$, $\alpha_e=2.2$):
\begin{center}
\begin{tabular}{lccc}
\toprule
 & $W$ [eV/ion] & efficiency $E_{\rm th}/W$ & $\langle N_{\rm ion}\rangle$ \\
\midrule
photons   & $21.075$ & $64.5\%$ & $1.0325$ \\
electrons & $48.038$ & $28.3\%$ & $122.92$ \\
\bottomrule
\end{tabular}
\end{center}
so $X_{\rm comp}=2.2794$: a photon buys an ionization for less than half the energy.

\paragraph{Where the electron's energy actually goes.} Averaged over the injection
spectrum, the loss flux divides as:
\begin{center}
\begin{tabular}{lr}
\toprule
mechanism & fraction of $dK/dt$ \\
\midrule
ionization & $72.07\%$ \\ 
excitation & $26.56\%$ \\ 
compton & $1.08\%$ \\ 
adiabatic & $0.22\%$ \\ 
coulomb & $0.07\%$ \\
\bottomrule
\end{tabular}
\end{center}
The decisive entry is \textbf{excitation at $26.56\%$}. That energy leaves as
Lyman-series photons, \emph{every one of which is below} $13.6$~eV --- Ly$\alpha$ is
$10.2$~eV --- so it can never ionize hydrogen, at any later time, by any route. It is not
delayed; it is gone. A further $1.37\%$ goes to inverse Compton on a CMB that is
$(1+z)^{4}\simeq1.5\times10^{4}$ times denser at $z=10$, plus adiabatic and Coulomb losses.
The remainder is sub-threshold kinetic energy: once a secondary electron drops below
$13.6$~eV it can only excite and heat, so it too leaves the ionization budget. Of the
energy injected, only $28.3\%$ ends up as binding energy delivered.

\paragraph{Why photons escape most of that penalty.} A photon does not avoid the electron
physics --- it \emph{defers} it. Photoionization converts the first $13.598$~eV at $100\%$
efficiency, and only the excess $E-E_{\rm th}$ is handed to a photoelectron, which from that
instant suffers exactly the same losses as any CR electron. Structurally,
\begin{equation}
  N_{{\rm ion},\gamma}(E)\;\simeq\;1+N_{{\rm ion},e}(E-E_{\rm th}),
  \label{eq:plusone}
\end{equation}
\emph{one free ionization plus an electron cascade}. Checked directly:
\begin{center}
\begin{tabular}{rrrr}
\toprule
$E$ [eV] & $N_{{\rm ion},\gamma}(E)$ & $1+N_{{\rm ion},e}(E-E_{\rm th})$ & rel.\ diff. \\
\midrule
54.4 & $1.858$ & $1.724$ & $7.8\%$ \\ 
100 & $2.963$ & $3.077$ & $3.7\%$ \\ 
1000 & $26.364$ & $28.977$ & $9.0\%$ \\ 
$10^{4}$ & $265.337$ & $277.162$ & $4.3\%$ \\ 
$10^{5}$ & $2157.676$ & $2203.241$ & $2.1\%$ \\
\bottomrule
\end{tabular}
\end{center}
Equation~\eqref{eq:plusone} is \emph{approximate}, not an identity: the few-per-cent
residual is the known disagreement between the depfit-route photon route and the loss+IC-route electron
route (assumption A9, $7.7\%$), not new physics.

\paragraph{The consequence.} The photon advantage is not that photons are intrinsically
efficient --- it is that they can be delivered \emph{at the threshold, where the job is
cheapest}, whereas CR electrons are injected at keV--TeV and must bleed down through
non-ionizing channels to get there. Monoenergetically the two costs converge above
$\sim1$~keV ($37.9$ vs $35.3$ eV/ion at $1$~keV; $46.3$ vs $45.4$ at $100$~keV), because a
keV photon's photoelectron \emph{is} a keV electron. The entire photon advantage lives in
the $13.6$--$100$~eV window, and is really a statement about how badly slow electrons
perform there: a $20$~eV photon ionizes with near-certainty, while a $20$~eV electron costs
$171$~eV per ionization because it overwhelmingly excites instead.

\section{The weakest input: $E_{\min,e}$}\label{sec:emin}
The two bands are not equally well determined. The photon band's floor is atomic physics and
its ceiling is derived from $\tau=1$; \emph{every} entry of the electron band is a choice
with no external source. That would be harmless if the answer were insensitive to it:
\begin{center}
\begin{tabular}{rrr}
\toprule
$E_{\min,e}$ [eV] & $W_e$ [eV/ion] & $X_{\rm comp}$ at fiducial indices \\
\midrule
$10^{2}$ & $43.066$ & $1.8585$ \\ 
$10^{3}$ & $48.038$ & $2.0730$ \\ 
$10^{4}$ & $60.929$ & $2.6293$ \\ 
$10^{5}$ & $94.470$ & $4.0768$ \\
\bottomrule
\end{tabular}
\end{center}
Lowering $E_{\min,e}$ drags in the catastrophically expensive sub-$100$~eV region; raising it
does the reverse. Across $10^{2}$--$10^{5}$~eV the whole map slides by a factor $2.19$
--- \textbf{more than the entire $\alpha_e$ axis moves it}. The map's \emph{shape} is well
posed; its \emph{normalisation} is not. The figure is therefore drawn at two values of
$E_{\min,e}$ rather than one, so the shift is visible rather than hidden in a constant.

\section{Inventory}
Status: \textbf{V} verified against a source in \texttt{papers/}; C cited, unverified;
D derived here; S scanned; U your choice; \textcolor{red}{\textbf{X}} unsourced.
\begin{center}
\begin{longtable}{lllll}
\toprule
symbol & value & units & depends on & st. \\
\midrule
\endhead
\multicolumn{5}{l}{\textbf{Scenario}} \\[2pt] 
$z$ & 10 & -- & -- & U \\ 
\multicolumn{5}{p{0.93\textwidth}}{\footnotesize\quad redshift of the snapshot\quad[src: project fiducial]} \\[3pt] 
$n_{\rm H}$ & $2.5287\times10^{-4}$ & cm$^{-3}$ & z, $\Omega_b h^2$, $X_{\rm H}$ & D \\ 
\multicolumn{5}{p{0.93\textwidth}}{\footnotesize\quad proper hydrogen number density\quad[src: Planck2020]} \\[3pt] 
$x_e$ & $10^{-4}$ & -- & -- & U \\ 
\multicolumn{5}{p{0.93\textwidth}}{\footnotesize\quad ionized fraction, held static\quad[src: your value]} \\[3pt] 
$H(z),\ t_H$ & 0.7086 Gyr & s$^{-1}$, Gyr & Planck18 complete $H(z)$ & D \\ 
\multicolumn{5}{p{0.93\textwidth}}{\footnotesize\quad expansion rate and Hubble time\quad[src: Planck2020]} \\[3pt] 
target & Z = 0 & -- & -- & U \\ 
\multicolumn{5}{p{0.93\textwidth}}{\footnotesize\quad pure atomic hydrogen at mean density\quad[src: project fiducial]} \\[3pt] 
\multicolumn{5}{l}{\textbf{Photon band}} \\[2pt] 
$E_{\min,\gamma}$ & 13.598434600 & eV & atomic constant & \textbf{V} \\ 
\multicolumn{5}{p{0.93\textwidth}}{\footnotesize\quad HI ionization potential -- PHYSICAL floor\quad[src: NIST\_ASD]} \\[3pt] 
E\_max,g (stellar) & 54.4228 & eV & band policy & D \\ 
\multicolumn{5}{p{0.93\textwidth}}{\footnotesize\quad band top ACTUALLY USED: min(4 Ryd source cutoff, IGM cutoff)\quad[src: DERIVED here]} \\[3pt] 
E\_max,g (IGM) & 1218.0 & eV & $\sigma_{\rm photoion}$, $n_{\rm HI}$, $c/H(z)$ & D \\ 
\multicolumn{5}{p{0.93\textwidth}}{\footnotesize\quad free-streaming onset: $\tau_{\rm IGM}$ = 1 over one Hubble length; caps any source whose own emission reaches past it\quad[src: DERIVED here]} \\[3pt] 
$\tau_{\rm IGM}(E)$ & 1.934 at 1 keV, 0.0476 at 3 keV & -- & $\sigma_{\rm photoion}$, $n_{\rm HI}$, z & D \\ 
\multicolumn{5}{p{0.93\textwidth}}{\footnotesize\quad IGM optical depth over a Hubble length\quad[src: Karzas1961 / DERIVED]} \\[3pt] 
$S_\gamma(E)$ & power law & eV$^{-1}$ & $\alpha_\gamma$, band & U \\ 
\multicolumn{5}{p{0.93\textwidth}}{\footnotesize\quad $dN/dE\propto E^{-(\alpha_\gamma+1)}$, NUMBER-normalised to 1 on the band\quad[src: convention]} \\[3pt] 
$\alpha_\gamma$ (fiducial) & 2.0 & -- & -- & U \\ 
\multicolumn{5}{p{0.93\textwidth}}{\footnotesize\quad the index actually used, $f_\nu\propto\nu^{-\alpha_\gamma}$. OUR PARAMETRISATION: MarquesChaves2026 gives the population (BPASS v2.2.1, Z=0.003, 6.8 Myr) but no LyC index, and its beta\_UV is measured longward of 1400 A, which cannot be carried across the Lyman break\quad[src: this work]} \\[3pt] 
$\alpha_\gamma$ (measured) & 4.5 +/- 0.4 & -- & -- & \textbf{V} \\ 
\multicolumn{5}{p{0.93\textwidth}}{\footnotesize\quad the ONLY direct measurement in our band: L\_nu ~ nu\^{}-gamma over 13.6-54.4 eV, exactly this project's convention, no conversion. CAVEAT: four GALACTIC HII regions at roughly solar metallicity, and the authors find it disagrees with BPASS; metal-poor z~10 should be harder, i.e. smaller\quad[src: Murchikova2020]} \\[3pt] 
$\alpha_\gamma$ (axis) & 0.3-5.0 & -- & -- & S \\ 
\multicolumn{5}{p{0.93\textwidth}}{\footnotesize\quad range swept by the figure, chosen to contain every index sourced in this project AND the measured 4.5\quad[src: ]} \\[3pt] 
\multicolumn{5}{l}{\textbf{Electron band}} \\[2pt] 
$E_{\min,e}$ & $10^{3}$ & eV & -- & U \\ 
\multicolumn{5}{p{0.93\textwidth}}{\footnotesize\quad low cutoff of the CR injection spectrum\quad[src: your choice]} \\[3pt] 
$E_{\max,e}$ & $10^{12}$ & eV & -- & U \\ 
\multicolumn{5}{p{0.93\textwidth}}{\footnotesize\quad high cutoff of the CR injection spectrum\quad[src: your choice]} \\[3pt] 
$S_e(E)$ & power law & eV$^{-1}$ & $\alpha_e$, band & U \\ 
\multicolumn{5}{p{0.93\textwidth}}{\footnotesize\quad $dN/dE\propto E^{-\alpha_e}$, NUMBER-normalised to 1 on the band\quad[src: convention]} \\[3pt] 
$\alpha_e$ & 2.2 (fid); 1.2-2.7 (axis) & -- & -- & U \\ 
\multicolumn{5}{p{0.93\textwidth}}{\footnotesize\quad CR electron injection index\quad[src: your choice; DSA predicts 2]} \\[3pt] 
\multicolumn{5}{l}{\textbf{Ionization yields}} \\[2pt] 
$N_{{\rm ion},\gamma}(E)$ & 1.0325 (band mean) & -- & E, $x_e$, z & \textbf{V} \\ 
\multicolumn{5}{p{0.93\textwidth}}{\footnotesize\quad ionizations per ABSORBED photon of energy E, primary + full secondary cascade\quad[src: \texttt{ionization\_yield}, depfit route]} \\[3pt] 
$N_{{\rm ion},e}(E)$ & 122.92 (band mean) & -- & E, $x_e$, z, B & \textbf{V} \\ 
\multicolumn{5}{p{0.93\textwidth}}{\footnotesize\quad ionizations per INJECTED electron of energy E, 7 loss channels + BED cascade + IC secondaries\quad[src: \texttt{igm\_losses} the loss+IC route]} \\[3pt] 
\multicolumn{5}{l}{\textbf{Derived moments}} \\[2pt] 
$\langle E\rangle_\gamma$ & 21.7598 & eV & $\alpha_\gamma$, band & D \\ 
\multicolumn{5}{p{0.93\textwidth}}{\footnotesize\quad mean energy per emitted ionizing photon\quad[src: DERIVED]} \\[3pt] 
$\langle E\rangle_e$ & $5.9049\times10^{3}$ & eV & $\alpha_e$, band & D \\ 
\multicolumn{5}{p{0.93\textwidth}}{\footnotesize\quad mean energy per injected CR electron\quad[src: DERIVED]} \\[3pt] 
$W_\gamma$ & 21.0747 & eV\,ion$^{-1}$ & $\alpha_\gamma$, band, $x_e$, z & D \\ 
\multicolumn{5}{p{0.93\textwidth}}{\footnotesize\quad energy price of one ionization, photons = $\langle E\rangle_\gamma/\langle N\rangle_\gamma$\quad[src: DERIVED]} \\[3pt] 
$W_e$ & 48.0378 & eV\,ion$^{-1}$ & $\alpha_e$, band, $x_e$, z, B & D \\ 
\multicolumn{5}{p{0.93\textwidth}}{\footnotesize\quad energy price of one ionization, electrons = $\langle E\rangle_e/\langle N\rangle_e$\quad[src: DERIVED]} \\[3pt] 
$X_{\rm comp}$ & 2.2794 & -- & $\alpha_\gamma$, $\alpha_e$ ONLY & D \\ 
\multicolumn{5}{p{0.93\textwidth}}{\footnotesize\quad $W_e/W_\gamma$ at $L_\gamma=L_e$: ionizations per unit energy, photons relative to electrons\quad[src: DERIVED]} \\[3pt]
\bottomrule
\end{longtable}
\end{center}
\end{document}
